% RETAINED REJECTED PROOF ATTEMPT. NOT INPUT INTO THE PAPER.
% The original Proved heading was provisional and incorrect as a proof claim.
% Independent referee boundary_d found that the proposed free-product formula
% ignores triple-cell attaching paths around a conductor cycle. Such a loop
% can map nontrivially to the Inoue component group. Root confirmed the flaw.
% The entire passage below is retained verbatim except this warning/heading.
% The replacement proof uses a hyperbolic link in the NEARBY fibre and the
% nonzero period of a relative holomorphic two-form, not this group formula.
\begin{proposition}[Rejected proof attempt: maximal monodromy with trivial central dualizing bundle]
Let $p:\mathcal X\to\Delta$ be a proper holomorphic map
from a smooth complex threefold, smooth over $\Delta^*$,
whose nonzero fibres are complex two-tori. Suppose that
$D=p^{-1}(0)$ is a reduced normal-crossing divisor and
\emph{assume in addition} that $\omega_D\simeq\OO_D$.
If its monodromy on $H_1(\mathcal X_t,\Q)$ satisfies
\[
 \operatorname{rank}(T-1)=2,                         \tag{B.4.28.1}
\]
then every normalization component of $D$ is a smooth
rational toric surface, its conductor is its toric
boundary, and the dual complex of $D$ is a topological
two-torus. In particular all normalized components are
K\"ahler.

No K\"ahler hypothesis on $\mathcal X$ is used. The
additional triviality of $\omega_D$ is not inferred from
the neighbouring two-tori. The proposition does not
produce a modification with this property from an
arbitrary analytic completion.
\end{proposition}
\begin{proof}
Write $V_i$ for the compact smooth normalization
components and $C_i$ for their reduced conductor divisors.
Adjunction gives
\[
 K_{V_i}+C_i\sim0.                                  \tag{B.4.28.2}
\]
If $D$ is smooth, the map is a proper submersion near
its compact central fibre and its monodromy is trivial.
Otherwise every $C_i$ is nonempty, by connectedness.
Adjunction on a conductor branch says that it is either
a smooth elliptic curve disjoint from the other branches,
or a rational curve with two boundary points, counted
on its normalization. The rational branches form cycles,
including the nodal irreducible and two-component cases.

\emph{Nonnegative Euler contributions.}
A nonempty effective anticanonical divisor makes every
pluricanonical system empty: multiplying a nonzero
pluricanonical section by the appropriate power of the
anticanonical section would give a nonzero holomorphic
function with a nonempty zero divisor. Thus the surface
has Kodaira dimension $-\infty$. Its minimal model is
rational, ruled, or of class VII. We spell out the
consequences for the reduced anticanonical pairs here.

On a rational surface, the conductor is connected, by
$0\to K_V\to\OO_V\to\OO_C\to0$ and $H^1(K_V)=0$.
For a rational cycle, the charge identity and positivity
are
\[
 e(V\setminus C)=12-C^2-\#\{\text{cycle components}\}
 \ge0,\qquad
 e(V\setminus C)=0\ \Longleftrightarrow\ (V,C)
 \text{ is toric}.                                  \tag{B.4.28.3}
\]
Here $e$ denotes compactly supported Euler characteristic.
This includes the nodal irreducible convention:
\citet[Definition 1.1 and Lemmas 1.2, 2.2, 2.7]{Friedman2016AB}.
For a smooth elliptic anticanonical curve,
$e(V\setminus C)=2+b_2(V)>0$ instead.

On a nonrational ruled minimal model, a conductor
component dominating the base has degree at most two
and genus at most one. The base is therefore elliptic.
The minimal conductor is two disjoint elliptic sections
or one unramified elliptic bisection. Its Euler
contribution is zero. The two section squares sum to
zero, and the bisection has square zero, also directly
from $K_V+C\sim0$ and $K_V^2=0$.

A minimal class-VII surface with $b_2=0$ and a curve is
a Hopf surface. Its elliptic curves have square zero,
and a reduced anticanonical divisor has at most two
elliptic components. Indeed its pullback to the finite
primary Hopf cover is a reduced anticanonical divisor,
the two elliptic axes. The Euler contribution is zero.
This permits a finite quotient interchanging the axes.

For a minimal class-VII surface with $b_2>0$, the
reduced anticanonical divisor is a numerically
anticanonical divisor whose coefficients are all one.
\citet[Corollary 2.25, Lemma 2.26 and Proposition
2.27]{Dloussky2005AB} imply that it is parabolic Inoue
or Inoue--Hirzebruch. Half-Inoue has only a nontrivially
twisted anticanonical divisor and cannot satisfy
(B.4.28.2). The remaining possibilities are hyperbolic
Inoue with two disjoint rational cycles, or parabolic
Inoue with one rational cycle and a disjoint elliptic
curve. If $n=b_2>0$, there are $n$ rational cycle
components and $e(V)=n$, so the contribution is zero.
In the parabolic case the elliptic curve has square
$-n$. The curve configurations, anticanonical identities
and underlying $(S^1\times S^3)\#n\overline{\PP}^2$
topology are recorded in
\citet[Sections 3.2--3.3]{FujikiPontecorvo2010AB}.

These assertions extend to nonminimal pairs. An
exceptional curve outside $C$ meets it once, transversely
at a smooth point, since $-K_V\cdot E=1$; its contraction
reverses an interior blow-up. An exceptional curve in
$C$ meets the rest at two points and its contraction
reverses a corner blow-up. The reduced anticanonical
condition is preserved. Interior blow-up increases
$e(V\setminus C)$ by one, whereas corner blow-up leaves
it unchanged. Thus every contribution is nonnegative.
Equality permits only toric rational pairs, minimal
elliptic ruled or Hopf pairs, and corner blow-ups of
hyperbolic or parabolic Inoue pairs. For these Inoue
pairs the cycle count and $b_2$ increase together, and
the fundamental group stays $\Z$. Corner blow-ups do
not touch the parabolic elliptic curve, whose square
remains strictly negative.

The local nearby fibre at a multiplicity-$k$ stratum
of a reduced normal-crossing divisor retracts to
$(S^1)^{k-1}$. Euler integration consequently gives
\[
 0=e(\mathcal X_t)=\sum_i e(V_i\setminus C_i).          \tag{B.4.28.4}
\]
Every summand is nonnegative and therefore zero.
Rational components with a smooth elliptic conductor,
and all interior blow-ups, are excluded.

\emph{Elliptic normal-degree balance.}
An elliptic double curve contains no triple point.
The local model $p=xy$ gives opposite degrees to its
two normal lines. Sum these degrees over all elliptic
conductor branches. The total is zero. An elliptic
ruled or Hopf component contributes zero; a parabolic
Inoue component contributes the strictly negative
degree of its sole elliptic branch. No rational
component with elliptic conductor remains. Hence there
are no parabolic Inoue components.

By connectedness, all conductors are now elliptic or
all are rational: no remaining component has both
types. In the elliptic case there are no triple points
and every vertex of the dual graph has valence at most
two. The graph is a segment or circle, allowing a loop
or multiple edges. In this case
\[
 \operatorname{rank}(T-1)\le b_1(\text{dual graph})
 \le1.                                               \tag{B.4.28.5}
\]
To see this without an invariant-cycle theorem, remove
small tubes about the disjoint double curves. The
remaining compact cores extend as smooth submersions
over the whole disc and can be trivialized there.
Monodromy is then a product of the annular twists
$xy=t$. A passage of a loop across one neck adds its
oriented normal-circle class, with coefficient the
signed crossing number of that edge. Core loops have
zero change. Crossing numbers of a closed loop form a
cycle of the dual graph. Thus $T-1$ factors through its
rational cycle space. The normal-circle rotation exists
on the circle bundle also for a nontrivial normal line;
a product normal bundle is not assumed. This proves
(B.4.28.5) and excludes the elliptic case.

\emph{Hyperbolic components and the fundamental group.}
The remaining components are toric rational pairs or
corner blow-ups of hyperbolic Inoue pairs. Let $h$ be
the number of the latter and $K$ the finite dual complex.
Each edge has two incident faces, since a rational
double curve has two triple points. A rational vertex
has one circular link; a hyperbolic vertex has two
disjoint circular links. Split each of these $h$ vertices
in two. The result is a disjoint union of closed
connected surfaces $\Sigma_1,\ldots,\Sigma_c$.

These surfaces are orientable. A trivialization of
$\omega_D$ pulls back to logarithmic two-forms on $V_i$.
Their iterated residues at triple points are nonzero.
On each rational double curve, the residues at its
two marked points are opposite, by the residue theorem
on $\PP^1$; gluing dualizing forms supplies the
opposite-branch sign. These coefficients give a
nowhere-zero two-cycle on every split surface, which
proves orientability. Let $g$ be the sum of their genera.

Coherent Euler characteristic is constant in a proper
flat analytic family. The alternating normalization
resolution gives
\[
 0=\chi(\OO_D)=\chi(K)-h.                             \tag{B.4.28.6}
\]
A rational surface contributes one, a hyperbolic Inoue
surface zero, a rational double curve one and a triple
point one with alternating signs. Splitting $h$ vertices
increases Euler characteristic by $h$, and hence
\[
 2c-2g=\chi(K)+h=2h,\qquad c-g=h.                     \tag{B.4.28.7}
\]
Recovering connected $K$ identifies $h$ pairs of points
in the disjoint surfaces. Thus
\[
 b_1(K)=2g+h-c+1=g+1\ge1.                            \tag{B.4.28.8}
\]
Van Kampen, with the incidence cylinders and triple-point
two-cells, gives
\[
 \pi_1(D)\simeq
 \bigl(\mathop{*}_{i=1}^{h}\Z\bigr)*\pi_1(K).           \tag{B.4.28.9}
\]
Indeed all double curves are simply connected. After
collapsing a maximal incidence tree, triple-point
relations involve only incidence generators; no
component group has an edge-group relation.

B.4.2 supplies a surjection
$\pi_1(\mathcal X_t)\twoheadrightarrow\pi_1(D)$.
The source is $\Z^4$, so the target is abelian.
If $h>0$, however, (B.4.28.8)--(B.4.28.9) give a
surjection from it to $\Z*\Z$, a contradiction.
Therefore $h=0$. All normalization pairs are toric,
the dual complex was already a connected closed
oriented surface, and (B.4.28.7) gives genus one.

For a self-crossing centre one uses the ordered local
branch incidence spaces and their cylinder realization.
A self-double rational curve gives an incidence edge
with trivial edge group. Subdividing the finite dual
cell structure displays the same circular links.
The alternating normalization resolution has the
branch signs and the same Euler count; no division by
a branch-permutation order is used. Thus the argument
also covers the stated normal-crossing scope.
\end{proof}

\begin{remark}[The missing reduction for an arbitrary completion]
The proposition cannot alone complete B.4.19 at its
missing origin. Suppose a relative holomorphic two-form
on a hypothetical transverse punctured slice has first
been proved to extend meromorphically. Rescale it so
that its vertical divisor is $\sum_i a_iV_i$ with
$\min_i a_i=0$. On each component its restriction gives
\[
 K_{V_i}+C_i\sim\sum_{j\ne i}(a_j-a_i)C_{ij}.          \tag{B.4.28.10}
\]
Here $\OO_{\mathcal X}(V_i)|_{V_i}\sim
-\sum_{j\ne i}C_{ij}$ since the full fibre is principal.
The minimal coefficient does not make this right side
zero: neighbouring higher-coefficient components remain.
Contracting them, while retaining semistability and
properness in an arbitrary non-K\"ahler threefold, is
not proved here. Nor is arbitrary meromorphic extension
of a punctured comparison assumed. The K\"ahler,
$K$-trivial torus classification of
\citet[Theorem 3.1]{KamenovaToriRoot} does not supply
that missing reduction.
\end{remark}

